\begin{section}{Fields and Ordered Fields}

A set $F$ is called a \textbf{field} if it is equipped with two operations,
\emph{addition} and \emph{multiplication}, and satisfies the following axioms:
\begin{enumerate}[label=\textbf{(\Alph*)}]
    \item Axioms for addition
    \begin{itemize}
        \item $x + y \in F$ for all $x, y \in F$.
        \item $x + y = y + x$ for all $x, y \in F$.
        \item $(x + y) + z = x + (y + z)$ for all $x, y, z \in F$.
        \item There exists an element $0 \in F$ such that $x + 0 = x$ for all $x \in F$.
        \item For each $x \in F$, there exists an element $-x \in F$ such that $x + (-x) = 0$.
    \end{itemize}
    \item Axioms for multiplication
    \begin{itemize}
        \item $x \cdot y \in F$ for all $x, y \in F$.
        \item $x \cdot y = y \cdot x$ for all $x, y \in F$.
        \item $(x \cdot y) \cdot z = x \cdot (y \cdot z)$ for all $x, y, z \in F$.
        \item There exists an element $1 \in F$ with $1 \neq 0$ such that $x \cdot 1 = x$ for all $x \in F$.
        \item For each $x \in F$ with $x \neq 0$, there exists an element $x^{-1} \in F$ such that $x \cdot x^{-1} = 1$.
    \end{itemize}
    \item The distributive law
    \begin{itemize}
        \item $x \cdot (y + z) = x \cdot y + x \cdot z$ for all $x, y, z \in F$.
    \end{itemize}
\end{enumerate}

A set $F$ is called an \textbf{ordered field} if it is both an ordered set and a field, and satisfies the following axioms:
\begin{itemize}
    \item If $x, y, z \in F$ and $x < y$, then $x + z < y + z$.
    \item If $x, y \in F$ and $x > 0$ and $y > 0$, then $x \cdot y > 0$.
\end{itemize}

\bigskip

\begin{example}
    $(\Q, +, \cdot)$ is an ordered field, while $(\Z, +, \cdot)$ is not an ordered field since $2^{-1} \notin \Z$.
\end{example}

\end{section}